
\documentclass[11pt]{article}
\usepackage[a4paper,margin=1in]{geometry}
\usepackage[T1]{fontenc}
\usepackage[utf8]{inputenc}
\usepackage{lmodern}
\usepackage{amsmath,amssymb,amsthm,mathtools}
\usepackage{microtype}
\usepackage{fancyhdr}
\usepackage{array,longtable,enumitem}
\setlength{\parindent}{1.5em}
\setlength{\parskip}{0.18em}
\emergencystretch=4em
\pagestyle{fancy}
\fancyhf{}
\cfoot{\thepage}

\usepackage[ngerman]{babel}
\lhead{Heinrich Weber, Lehrbuch der Algebra, Band III}
\rhead{deutsche Abschrift}
\begin{document}
\begin{center}
{\large Erster Abschnitt.}\\[0.7em]
{\Large Die elliptischen Integrale.}
\end{center}

\section*{\S\,1. Definition der elliptischen Integrale.}

Wenn die systematische Darstellung der Integralrechnung bis zu dem Punkte gelangt ist, wo algebraische Integrale mit der Quadratwurzel aus einer Funktion ersten oder zweiten Grades auf Integrale rationaler Funktionen zurückgeführt werden, so tritt an dieser Stelle dem Lernenden eine Schranke entgegen, die er mit den ihm bis dahin zu Gebote stehenden Hilfsmitteln nicht zu übersteigen imstande ist. Das Streben nach einer Erweiterung der Hilfsmittel, um auch noch die nächste Klasse von Integralen der Forschung zugänglich zu machen, ist, wie es historisch der Anlaß gewesen, zu einem eingehenderen Studium der elliptischen Integrale und zur Einführung der elliptischen Funktionen, auch der naturgemäßeste und verständlichste Ausgangspunkt für den, der in die Theorie dieser Funktionen zuerst eingeführt werden soll. Es soll daher auch unsere nächste Aufgabe sein, uns mit den elliptischen Integralen und ihren wichtigsten Eigenschaften bekannt zu machen.

Die Definition eines elliptischen Integrals, von der wir ausgehen wollen, ist die folgende: Es bedeute $f(x)$ eine ganze rationale Funktion dritten oder vierten Grades mit vier verschiedenen Wurzeln
\begin{equation}
f(x)=a_0x^4+4a_1x^3+6a_2x^2+4a_3x+a_4,
\tag{1}
\end{equation}
worin $a_0$ und $a_1$ nicht beide zugleich verschwinden; es sei ferner $\Phi(x,y)$ eine beliebige ganze oder gebrochene rationale Funktion der beiden Argumente $x,y$. Dann ist
\[
\int \Phi(x,\sqrt{f(x)})\,dx
\]
das allgemeine elliptische Integral und
\[
\Phi(x,\sqrt{f(x)})\,dx
\]
das allgemeine elliptische Differential. Es läßt sich nun aber dies allgemeine Integral und Differential auf wesentlich einfachere zurückführen. Zunächst läßt sich $\Phi$ in die Form setzen:
\[
\frac{A+B\sqrt{f(x)}}{C+D\sqrt{f(x)}},
\]
worin $A,B,C,D$ ganze rationale Funktionen von $x$ sind, und indem man diesen Bruch mit $C-D\sqrt{f(x)}$ erweitert, in die Form
\[
\Psi(x)+\frac{\Phi(x)}{\sqrt{f(x)}},
\]
worin $\Psi(x)$ und $\Phi(x)$ ganze oder gebrochene rationale Funktionen von $x$ sind, nämlich
\[
\Psi(x)=\frac{AC-BD f(x)}{C^2-D^2 f(x)},
\]
\[
\Phi(x)=\frac{(BC-AD)f(x)}{C^2-D^2 f(x)}.
\]

Wir lassen nun das rationale Integral
\[
\int \Psi(x)\,dx,
\]
das auf Logarithmen und algebraische Funktionen führt, außer Betracht, und befassen uns nur noch mit dem elliptischen Integral
\begin{equation}
\int \frac{\Phi(x)\,dx}{\sqrt{f(x)}}
\tag{2}
\end{equation}
oder dem entsprechenden elliptischen Differential
\begin{equation}
\Phi(x)\frac{dx}{\sqrt{f(x)}}.
\tag{3}
\end{equation}

Es ist bisweilen nützlich, dies Differential in der homogenen Form zu betrachten; zu diesem Ende setzen wir für $x$ das Verhältnis zweier Variablen $x:y$, und demnach für $dx$
\[
\frac{y\,dx-x\,dy}{y^2}.
\]
Ist dann
\[
f(x,y)=a_0x^4+4a_1x^3y+6a_2x^2y^2+4a_3xy^3+a_4y^4,
\]
und $\Phi(x,y)$ eine homogene Funktion $0$ter Ordnung, so erhält unser elliptisches Differential die Gestalt
\begin{equation}
\Phi(x,y)\frac{y\,dx-x\,dy}{\sqrt{f(x,y)}}.
\tag{4}
\end{equation}

Die Aufgabe, die uns nun zunächst beschäftigen wird, ist eine doppelte: Es soll durch Einführung neuer Veränderlicher das Differential
\begin{equation}
\frac{dx}{\sqrt{f(x)}}\quad\hbox{oder}\quad \frac{y\,dx-x\,dy}{\sqrt{f(x,y)}}
\tag{5}
\end{equation}
in ein anderes von derselben Form transformiert werden, worin aber die Funktion unter dem Quadratwurzelzeichen möglichst vereinfacht und besonders von einer möglichst kleinen Zahl von Parametern abhängig gemacht wird (\S\S~3, 4, 5). Es soll zweitens das allgemeine Differential (3) oder (4) in andere ähnliche zerlegt werden, in welchen die Funktion $\Phi(x)$ oder $\Phi(x,y)$ möglichst einfach ist (\S~12).

Die Funktion $f(x)$ läßt sich in lineare Faktoren zerlegen (Bd.~I, \S~33). Wir setzen demnach auch
\[
f(x)=a_0(x-x_1)(x-x_2)(x-x_3)(x-x_4).
\]
Denken wir uns die Größen $x,x_1,x_2,\ldots$ als Abszissen von Punkten auf eine gerade Linie von einem beliebigen Nullpunkte aus aufgetragen, die positiven nach der einen, die negativen nach der anderen Seite (z.~B. positiv nach rechts, negativ nach links), so stellen die Differenzen $x-x_1$, $x-x_2$, $x-x_3$, $x-x_4$ die Entfernungen des Punktes $x$ von $x_1,x_2,x_3,x_4$ dar, und zwar positiv gerechnet, wenn $x$ nach der positiven Seite von $x_1,x_2,\ldots$ liegt. In demselben Sinne ist $x_2-x_1$ die Entfernung des Punktes $x_2$ von $x_1$ usf.

Wir bezeichnen die Differenzen auch kürzer durch $(xx_1)$, $(xx_2),\ldots,(x_2x_1),\ldots$. Wir benutzen diese Darstellungsweise aber nur zur leichteren Übersicht und schließen auch imaginäre Punkte der Geraden nicht aus.

\section*{\S\,2. Doppelverhältnisse.}

Das einfachste Mittel zur Vereinfachung des elliptischen Differentials stützt sich auf die lineare Transformation der Funktion $f(x)$ (Bd.~I, \S~67). Es mögen $\alpha,\beta,\gamma,\delta$ vier beliebige Konstanten bedeuten, deren Determinante
\begin{equation}
\alpha\delta-\beta\gamma=r
\tag{1}
\end{equation}
von Null verschieden ist. Wir setzen
\begin{equation}
x=\frac{\alpha x'+\beta}{\gamma x'+\delta},
\tag{2}
\end{equation}
und nennen dies eine lineare Substitution. Deuten wir $x$ und $x'$ wie im \S~1 als Punkte auf zwei geraden Linien $L$ und $L'$, so ist durch (2) eine gegenseitig eindeutige Beziehung der Punkte von $L$ und $L'$, eine Abbildung, festgelegt. Den Werten $x=\infty$ und $x'=\infty$ entsprechen die unendlich fernen Teile der Geraden, die wir ebenfalls als bestimmte Punkte betrachten. Es entspricht dann der Punkt $x'=\infty$ dem Punkt $x=\alpha:\gamma$ und der Punkt $x=\infty$ dem Punkt $x'=-\delta:\gamma$, und nur wenn $\gamma=0$, die Substitution (2) also ganz ist, entsprechen sich die unendlich fernen Punkte auf $L$ und $L'$ gegenseitig.

Die Substitution (2) ändert sich nicht, wenn die vier Transformationszahlen mit demselben Faktor multipliziert werden. Die Determinante $r$ vervielfältigt sich mit dem Quadrate dieses Faktors, und man kann daher diesen Faktor so bestimmen, daß die Determinante einen beliebig gegebenen Wert, z.~B. den Wert $1$ erhält.

Wendet man die Substitution (2) auf irgend zwei Punkte $x_1,x_2$ und die zugehörigen $x'_1,x'_2$ an, so ergibt sich
\[
(x_1x_2)=\frac{r(x'_1x'_2)}{(\gamma x'_1+\delta)(\gamma x'_2+\delta)},
\]
und wenn man ein zweites Punktepaar $x_3,x_4$ und das entsprechende $x'_3,x'_4$ hinzunimmt:
\[
(x_1x_2)(x_3x_4)=\frac{r^2(x'_1x'_2)(x'_3x'_4)}{(\gamma x'_1+\delta)(\gamma x'_2+\delta)(\gamma x'_3+\delta)(\gamma x'_4+\delta)}.
\]
Vertauscht man hierin $x_2$ mit $x_3$ und bildet den Quotienten, so folgt
\begin{equation}
\frac{(x_1x_2)(x_3x_4)}{(x_1x_3)(x_2x_4)}=\frac{(x'_1x'_2)(x'_3x'_4)}{(x'_1x'_3)(x'_2x'_4)}.
\tag{3}
\end{equation}
Der Ausdruck auf der linken Seite wird das Doppelverhältnis des Punktepaares $x_1x_4$ zu dem Punktepaar $x_2x_3$ genannt, und es ist also in dieser Formel der Satz enthalten:

\begin{center}
\textbf{Das Doppelverhältnis zweier Punktepaare ist bei gleichzeitiger linearer Transformation der vier Punkte invariant.}
\end{center}

Vier Punkte lassen sich auf sechs Arten in zwei Paare zerlegen. Setzen wir aber
\[
a=(x_2x_3)(x_1x_4),\qquad b=(x_3x_1)(x_2x_4),\qquad c=(x_1x_2)(x_3x_4),
\]
so besteht die Identität:
\begin{equation}
a+b+c=0
\tag{4}
\end{equation}
und wir erhalten die sechs Doppelverhältnisse
\begin{equation}
-\frac{c}{b},\quad -\frac{a}{c},\quad -\frac{b}{a},\quad -\frac{b}{c},\quad -\frac{c}{a},\quad -\frac{a}{b};
\tag{5}
\end{equation}
setzen wir das erste von ihnen $-c:b=\varkappa^2$, so erhält man nach (4) die sechs:
\begin{equation}
\varkappa^2,\quad 1-\frac1{\varkappa^2},\quad \frac1{1-\varkappa^2},\quad \frac1{\varkappa^2},\quad \frac{\varkappa^2}{\varkappa^2-1},\quad 1-\varkappa^2.
\tag{6}
\end{equation}
Es sind also die sechs Doppelverhältnisse aus den vier Punkten linear durch eines unter ihnen ausgedrückt.

Wenn die $x_1,x_2,x_3,x_4$ untereinander permutiert werden, so werden die $a,b,c$ untereinander permutiert und ändern ihre Vorzeichen, jedoch so, daß entweder alle drei Vorzeichen gleichzeitig geändert werden oder ungeändert bleiben. Beispielsweise geben die Transpositionen
\[
\begin{array}{c@{\quad}c}
(14)\text{ und }(23)&\begin{pmatrix}a&b&c\\-a&-c&-b\end{pmatrix}\\[0.7em]
(24)\text{ und }(31)&\begin{pmatrix}a&b&c\\-c&-b&-a\end{pmatrix}\\[0.7em]
(34)\text{ und }(12)&\begin{pmatrix}a&b&c\\-b&-a&-c\end{pmatrix}.
\end{array}
\]
Wenn $x_1,x_2,x_3,x_4$ reell sind, so wird $\varkappa^2$ dann und nur dann ein positiver echter Bruch, wenn $\varkappa^2$ und $1-\varkappa^2$ positiv, also entweder $a$ und $c$ positiv und $b$ negativ oder $a$ und $c$ negativ, $b$ positiv ist. Dies findet statt, wenn $x_1,x_2,x_3,x_4$ in dieser Reihenfolge der Größe nach aufsteigend einander folgen und bei allen den Anordnungen, die daraus durch solche Permutationen entstehen, die $b$ ungeändert lassen. Dies sind die folgenden acht:
\[
\begin{array}{cccc}
1&2&3&4\\
2&3&4&1\\
3&4&1&2\\
4&1&2&3\\
1&4&3&2\\
4&3&2&1\\
3&2&1&4\\
2&1&4&3
\end{array}
\]
Denkt man sich also $x_1,x_2,x_3,x_4$ in dieser Reihenfolge auf eine Kreisperipherie gesetzt, so erhält man immer dann ein positives echt gebrochenes $\varkappa^2$, wenn die $x_i$ so der Größe nach aufeinander folgen, daß man mit einem beliebigen als kleinstem anfängt und dann auf dem Kreise entweder nach rechts oder nach links weiter zählt. Wir drücken dies so aus:

\begin{center}
\textbf{Damit $\varkappa^2$ ein positiver echter Bruch sei, müssen die Größen $x_1,x_2,x_3,x_4$ der Größe nach zyklisch aufeinander folgen.}
\end{center}


\section*{\S\,3. Lineare Transformation des elliptischen Differentials.}

Wenn sich die beiden Punkte $x_3$ und $x_4$ einem und demselben Punkte $x$ annähern, so nähern sich $x'_3$ und $x'_4$ dem entsprechenden Punkte $x'$ an. Wir setzen dann $(x_3x_4):(x'_3x'_4)=dx:dx'$ und erhalten aus (3), \S~2, die Beziehung zwischen den Differentialen $dx,dx'$:
\begin{equation}
\frac{(x_1x_2)\,dx}{(x_1x)(x_2x)}=\frac{(x'_1x'_2)\,dx'}{(x'_1x')(x'_2x')}.
\tag{1}
\end{equation}
Ebenso ergibt sich, wenn man an Stelle von $x_1,x_2$ zwei andere Punkte $x_3,x_4$ setzt:
\begin{equation}
\frac{(x_3x_4)\,dx}{(x_3x)(x_4x)}=\frac{(x'_3x'_4)\,dx'}{(x'_3x')(x'_4x')},
\tag{2}
\end{equation}
und wenn man multipliziert und die Wurzel zieht:
\begin{equation}
\frac{\sqrt{(x_1x_2)(x_3x_4)}\,dx}{\sqrt{(x_1x)(x_2x)(x_3x)(x_4x)}}=
\frac{\sqrt{(x'_1x'_2)(x'_3x'_4)}\,dx'}{\sqrt{(x'_1x')(x'_2x')(x'_3x')(x'_4x')}} ,
\tag{3}
\end{equation}
also eine Transformation des elliptischen Differentials durch eine lineare Substitution:
\[
x=\frac{\alpha x'+\beta}{\gamma x'+\delta}.
\]

Diese Substitution ist vollkommen bestimmt, wenn zu drei beliebigen Punkten $x_1,x_2,x_3$ die zugehörigen Werte $x'_1,x'_2,x'_3$ willkürlich gegeben sind, und man erhält sie aus der Gleichheit des Doppelverhältnisses:
\begin{equation}
\frac{(x_1x_2)(x_3x)}{(x_1x_3)(x_2x)}=\frac{(x'_1x'_2)(x'_3x')}{(x'_1x'_3)(x'_2x')},
\tag{4}
\end{equation}
durch Auflösung nach $x$, und den vierten zu $x_4$ gehörigen Wert $x'_4$ erhält man aus
\begin{equation}
\frac{(x_1x_2)(x_3x_4)}{(x_1x_3)(x_2x_4)}=\frac{(x'_1x'_2)(x'_3x'_4)}{(x'_1x'_3)(x'_2x'_4)}.
\tag{5}
\end{equation}

Wenn die Funktion $f(x)$ in dem elliptischen Differential
\[
du=\frac{dx}{\sqrt{f(x)}}
\]
gegeben ist, so sind damit auch die $x_1,x_2,x_3,x_4$ gegeben, und durch die lineare Transformation (3) kann man das Differential auf eine andere Form bringen, bei der drei der Größe $x'_1,x'_2,x'_3,x'_4$ beliebig gegebene Werte haben. Man erhält die Normalform, wenn man
\[
x'=z,\qquad x'_1=0,\qquad x'_2=1,\qquad x'_3=\frac{1}{\varkappa^2},\qquad x'_4=\infty
\]
annimmt. Die Größe $\varkappa$ heißt bei Legendre der Modul des elliptischen Differentials und $\sqrt{1-\varkappa^2}=\varkappa'$ das Komplement des Moduls. Ordnet man also die Punkte in folgender Weise einander zu:
\[
\begin{array}{ccccc}
z,&0,&1,&\dfrac{1}{\varkappa^2},&\infty,\\[0.4em]
x,&x_1,&x_2,&x_3,&x_4,
\end{array}
\]
so ergeben sich aus der Gleichheit der Doppelverhältnisse leicht die Relationen:
\begin{align}
z&=\frac{(xx_1)(x_2x_4)}{(xx_4)(x_2x_1)},\notag\\
1-z&=\frac{(xx_2)(x_1x_4)}{(xx_4)(x_1x_2)},\tag{6}\\
1-\varkappa^2z&=\frac{(xx_3)(x_1x_4)}{(xx_4)(x_1x_3)},\notag
\end{align}
\begin{equation}
dz=\frac{(x_4x_1)(x_4x_2)}{(x_2x_1)}\,\frac{dx}{(xx_4)^2},
\tag{7}
\end{equation}
\begin{equation}
\varkappa^2=\frac{(x_3x_4)(x_1x_2)}{(x_1x_3)(x_2x_4)},\qquad
\varkappa'^2=\frac{(x_2x_3)(x_1x_4)}{(x_2x_4)(x_1x_3)},
\tag{8}
\end{equation}
und aus (3):
\begin{equation}
\frac{\sqrt{(x_3x_1)(x_4x_2)}\,dx}{\sqrt{-(x_1x)(x_2x)(x_3x)(x_4x)}}=
\frac{dz}{\sqrt{z(1-z)(1-\varkappa^2z)}}.
\tag{9}
\end{equation}
Nach \S~2 wird $\varkappa^2$ ein positiver echter Bruch, wenn die $x_1,x_2,x_3,x_4$ reell sind und der Größe nach zyklisch aufeinander folgen. Dies gibt also acht verschiedene solche Transformationen, und man kann darunter je zwei auswählen, bei denen das Intervall $z=0$ bis $z=1$ einem gegebenen der Intervalle $(x_1x_2),(x_2x_3),(x_3x_4),(x_4x_1)$ entspricht; dem wachsenden $z$ entsprechen bei der einen dieser Transformationen die wachsenden $x$, bei der anderen die abnehmenden $x$ (mit dem Durchgange von $+\infty$ zu $-\infty$). Im ersten Falle haben die Quadratwurzeln in (9) beiderseits das gleiche, im zweiten das entgegengesetzte Zeichen.

Wenn man nicht darauf besteht, daß $\varkappa^2$ ein positiver echter Bruch ist, so kann man die Punkte $x_1,x_2,x_3,x_4$ auf alle Arten permutieren und erhält 24 verschiedene lineare Transformationen in die Normalform, von denen je vier dasselbe $\varkappa^2$ ergeben; man erhält im ganzen sechs verschiedene Moduln, die nach \S~2, (6) auseinander abgeleitet werden.

Man übersieht am leichtesten die Gesamtheit dieser Transformationen, wenn man annimmt, das zu transformierende Integral habe bereits die Normalform:
\[
\frac{dx}{\sqrt{x(1-x)(1-\lambda^2x)}};
\]
man hat dann in den Formeln (6) bis (9) die $x_1,x_2,x_3,x_4$ auf alle möglichen Arten durch $0,1,1:\lambda^2,\infty$ zu ersetzen. Wir nehmen als Beispiel die folgenden Zuordnungen:
\[
\begin{array}{c|cccc}
&x_1&x_2&x_3&x_4\\ \hline
1)&0&\infty&\dfrac{1}{\lambda^2}&1\\[0.5em]
2)&0&\dfrac{1}{\lambda^2}&1&\infty\\[0.5em]
3)&1&0&\infty&\dfrac{1}{\lambda^2}.
\end{array}
\]
Man erhält dann aus (6) und (8)
\[
\begin{array}{r@{\quad}l@{\quad}l}
1)& z=-\dfrac{x}{1-x},& \varkappa^2=1-\lambda^2=\lambda'^2,\\[0.7em]
2)& z=\lambda^2x,& \varkappa^2=\dfrac1{\lambda^2},\\[0.7em]
3)& z=\dfrac{1-x}{1-\lambda^2x},& \varkappa^2=\lambda^2,
\end{array}
\]
und für das Differential
\[
\frac{dz}{\sqrt{z(1-z)(1-\varkappa^2z)}}
\]
ergibt sich in den drei Fällen:
\[
\begin{array}{r@{\quad}c}
1)& \dfrac{dx}{\sqrt{-x(1-x)(1-\lambda^2x)}},\\[0.9em]
2)& \dfrac{\lambda\,dx}{\sqrt{x(1-x)(1-\lambda^2x)}},\\[0.9em]
3)& \dfrac{dx}{\sqrt{x(1-x)(1-\lambda^2x)}}.
\end{array}
\]

\section*{\S\,4. Die Legendresche Normalform.}

Wenn die Funktion $f(x)$, die in dem elliptischen Differential unter dem Wurzelzeichen steht, reelle Koeffizienten hat, so sind drei Fälle möglich:
\begin{enumerate}[label=\arabic*.]
\item $f(x)$ hat vier reelle Wurzeln $x_1,x_2,x_3,x_4$;
\item $f(x)$ hat zwei reelle Wurzeln $x_1,x_2$ und ein paar konjugiert imaginäre Wurzeln $x_3,x_4$;
\item $f(x)$ hat zwei Paare konjugiert imaginärer Wurzeln $x_1,x_2$ und $x_3,x_4$.
\end{enumerate}
In allen drei Fällen läßt sich das elliptische Differential durch eine reelle lineare Transformation auf die Form bringen:
\begin{equation}
\frac{dz}{\sqrt{(z^2-\alpha)(z^2-\beta)}},
\tag{1}
\end{equation}
worin $\alpha$ und $\beta$ reelle (positive oder negative) Konstanten sind.

Wir bestimmen die lineare Abhängigkeit zwischen den Variablen $x$ und $z$ so, daß sich folgende Werte entsprechen:
\begin{equation}
\begin{array}{ccccc}
x,&x_1,&x_2,&x_3,&x_4,\\
z,&\sqrt{\alpha},&-\sqrt{\alpha},&\sqrt{\beta},&-\sqrt{\beta},
\end{array}
\tag{2}
\end{equation}
und erhalten eine Substitution der Form:
\begin{equation}
h\frac{x-x_1}{x-x_2}=\frac{z-\sqrt{\alpha}}{z+\sqrt{\alpha}}.
\tag{3}
\end{equation}
Um $h$ zu bestimmen, setzen wir $x=x_3$, $x=x_4$ und entsprechend $z=\sqrt{\beta}$, $z=-\sqrt{\beta}$, und erhalten
\[
h\frac{(x_3x_1)}{(x_3x_2)}=\frac{\sqrt{\beta}-\sqrt{\alpha}}{\sqrt{\beta}+\sqrt{\alpha}},\qquad
h\frac{(x_4x_1)}{(x_4x_2)}=\frac{\sqrt{\beta}+\sqrt{\alpha}}{\sqrt{\beta}-\sqrt{\alpha}}.
\]
Daraus durch Multiplikation und Division
\begin{equation}
h=\sqrt{\frac{(x_3x_2)(x_4x_2)}{(x_3x_1)(x_4x_1)}},\qquad
\frac{\sqrt{\alpha}-\sqrt{\beta}}{\sqrt{\alpha}+\sqrt{\beta}}=
\sqrt{\frac{(x_3x_1)(x_4x_2)}{(x_3x_2)(x_4x_1)}}.
\tag{4}
\end{equation}
Setzen wir, wie in \S~2,
\begin{equation}
\begin{aligned}
a&=(x_2x_3)(x_1x_4),\\
b&=(x_3x_1)(x_2x_4),\\
c&=(x_1x_2)(x_3x_4),\\
a+b+c&=0,
\end{aligned}
\tag{5}
\end{equation}
so können wir, da es nur auf das Verhältnis von $\alpha$ zu $\beta$ ankommt, die letzte der Gleichungen (4) dadurch befriedigen, daß wir setzen:
\begin{equation}
\sqrt{\alpha}=\sqrt{\pm a}+\sqrt{\mp b},\qquad
\sqrt{\beta}=\sqrt{\pm a}-\sqrt{\mp b},
\tag{6}
\end{equation}
und aus (3) ergibt sich für $z$ der Ausdruck
\begin{equation}
z=\sqrt{\alpha}\,
\frac{(xx_2)\sqrt{(x_3x_1)(x_4x_1)}+(xx_1)\sqrt{(x_3x_2)(x_4x_2)}}
{(xx_2)\sqrt{(x_3x_1)(x_4x_1)}-(xx_1)\sqrt{(x_3x_2)(x_4x_2)}}.
\tag{7}
\end{equation}
Stellt man neben (3) noch die daraus folgende Gleichung:
\[
h'\frac{x-x_3}{x-x_4}=\frac{z-\sqrt{\beta}}{z+\sqrt{\beta}}
\]
auf, so ergibt sich durch logarithmische Differentiation:
\[
\frac{(x_1x_2)dx}{(xx_1)(xx_2)}=\frac{2\sqrt{\alpha}\,dz}{z^2-\alpha},\qquad
\frac{(x_3x_4)dx}{(xx_3)(xx_4)}=\frac{2\sqrt{\beta}\,dz}{z^2-\beta},
\]
und daraus durch Multiplikation mit Rücksicht auf (5) und (6):
\begin{equation}
\frac{dx}{\sqrt{\pm (xx_1)(xx_2)(xx_3)(xx_4)}}=
\frac{2dz}{\sqrt{(z^2-\alpha)(z^2-\beta)}}.
\tag{8}
\end{equation}

Wenn nun die vier Wurzeln von $f(x)$ reell sind, so gibt es, wie wir im \S~2 gesehen haben, acht Arten, diese Wurzeln den Zeichen $x_1,x_2,x_3,x_4$ so zuzuordnen, daß $a$ und $b$ entgegengesetzte Zeichen haben, und wenn also $\pm a,\mp b$ positiv sind, so werden $\sqrt{\alpha},\sqrt{\beta}$ reell, also $\alpha,\beta$ positiv, und nach (7) wird dann auch $z$ reell.

Sind zweitens $x_1,x_2$ reell, $x_3,x_4$ konjugiert imaginär, so sind $a$ und $-b$ konjugiert imaginär, also wird, wenn die Vorzeichen der Quadratwurzeln $\sqrt{\pm a},\sqrt{\mp b}$ passend bestimmt werden, $\sqrt{\alpha}$ reell, $\sqrt{\beta}$ rein imaginär, also $\alpha$ positiv, $\beta$ negativ und $z$ reell [weil $(x_3x_1)(x_4x_1)$ und $(x_3x_2)(x_4x_2)$ als Produkte konjugiert imaginärer Größen positiv sind].

Sind endlich $x_1,x_2$ und $x_3,x_4$ zwei konjugiert imaginäre Paare, so sind $a$ und $-b$ beide positiv. Nehmen wir also in (6) die unteren Zeichen, so werden $\sqrt{\alpha},\sqrt{\beta}$ rein imaginär, also $\alpha$ und $\beta$ negativ, und aus (7) ergibt sich für $z$ ein reeller Ausdruck.

Setzt man dann
\begin{equation}
z^2=y,
\tag{9}
\end{equation}
so ergibt sich
\begin{equation}
\frac{2dz}{\sqrt{(z^2-\alpha)(z^2-\beta)}}=
\frac{dy}{\sqrt{y(y-\alpha)(y-\beta)}}
\tag{10}
\end{equation}
und man hat also durch die quadratische Substitution (9) ein elliptisches Differential erhalten, bei dem unter dem Wurzelzeichen eine Funktion dritten Grades mit reellen Wurzeln steht, das man nach \S~3 durch lineare Substitution auf die Normalform
\[
\frac{d\xi}{\sqrt{\xi(1-\xi)(1-\varkappa^2\xi)}}
\]
bringen kann, und darin können $\xi$ und $\varkappa$ als positive echte Brüche angenommen werden.

Die Legendresche Normalform ergibt sich daraus, wenn man
\[
\xi=\sin^2\varphi,\qquad d\xi=2\sin\varphi\cos\varphi\,d\varphi
\]
setzt:
\begin{equation}
\frac{d\xi}{\sqrt{\xi(1-\xi)(1-\varkappa^2\xi)}}=
\frac{2d\varphi}{\sqrt{1-\varkappa^2\sin^2\varphi}}.
\tag{11}
\end{equation}

\section*{\S\,5. Die Weierstrasssche Normalform.}

Eine andere Normalform des elliptischen Differentials hat Weierstrass seinen Untersuchungen zugrunde gelegt, nämlich die Form
\begin{equation}
du=\frac{dz}{\sqrt{4z^3-g_2z-g_3}},
\tag{1}
\end{equation}
in der $g_2,g_3$ Konstanten sind, die die erste und zweite Invariante des Differentials genannt werden. Das allgemeine elliptische Differential
\begin{equation}
\frac{dx}{\sqrt{f(x)}},
\tag{2}
\end{equation}
worin
\begin{equation}
f(x)=a_0x^4+a_1x^3+a_2x^2+a_3x+a_4
\tag{3}
\end{equation}
ist, kann durch eine lineare Transformation auf die Weierstrasssche Normalform gebracht werden, wenn die Wurzeln $x_1,x_2,x_3,x_4$ von $f(x)$ bekannt sind. Am einfachsten geschieht dies auf folgende Weise.

Die Wurzeln der kubischen Funktion
\[
\varphi(z)=4z^3-g_2z-g_3
\]
mögen mit $e_1,e_2,e_3$ bezeichnet sein. Dann ist
\[
\varphi(z)=4(z-e_1)(z-e_2)(z-e_3)
\]
und $e_1+e_2+e_3=0$.

Wir lassen die Werte von $x$ und $z$ einander folgendermaßen entsprechen:
\[
\begin{array}{ccccc}
x,&x_1,&x_2,&x_3,&x_4,\\
z,&\infty,&e_1,&e_2,&e_3,
\end{array}
\]
und es ergibt sich, wenn wir mit $m$ einen konstanten Faktor bezeichnen, der willkürlich angenommen werden kann:
\begin{equation}
\begin{aligned}
z-e_1&=\frac{m(xx_2)}{(x_1x_2)(xx_1)}=m\left(\frac1{(x_1x_2)}-\frac1{(x_1x)}\right),\\
z-e_2&=\frac{m(xx_3)}{(x_1x_3)(xx_1)}=m\left(\frac1{(x_1x_3)}-\frac1{(x_1x)}\right),\\
z-e_3&=\frac{m(xx_4)}{(x_1x_4)(xx_1)}=m\left(\frac1{(x_1x_4)}-\frac1{(x_1x)}\right).
\end{aligned}
\tag{4}
\end{equation}
Der Faktor $m$ muß in allen drei Gleichungen derselbe sein, damit die Differenzen $(z-e_3)-(z-e_2)=e_2-e_3$ usw. von $x$ unabhängig werden.

Bildet man diese Differenzen und setzt wie in Algebra Bd.~I, \S~70
\[
(x_1x_2)(x_3x_4)=U,\qquad (x_1x_3)(x_4x_2)=V,\qquad (x_1x_4)(x_2x_3)=W,
\]
so folgt
\[
e_2-e_3=\frac{-mU}{(x_1x_2)(x_1x_3)(x_1x_4)}.
\]
Führt man einen neuen willkürlichen Faktor $\mu$ ein, indem man
\begin{equation}
m=3\mu a_0(x_1x_2)(x_1x_3)(x_1x_4)
\tag{5}
\end{equation}
setzt, so folgt
\[
e_2-e_3=-3\mu a_0U,\qquad e_3-e_1=-3\mu a_0V,\qquad e_1-e_2=-3\mu a_0W
\]
und daraus wegen $e_1+e_2+e_3=0$,
\[
e_1=\mu a_0(V-W),\qquad e_2=\mu a_0(W-U),\qquad e_3=\mu a_0(U-V).
\]
Die in Algebra I, \S~70 eingeführten Größen $y_1,y_2,y_3$ sind also gleich $e_1:\mu$, $e_2:\mu$, $e_3:\mu$, und man erhält nach der dortigen Formel (11) für die $e_1,e_2,e_3$ die kubische Gleichung:
\[
z^3-3A\mu^2z+\mu^3B=0.
\]
Es wird also, wenn $\varphi(z)=z^3-g_2z-g_3$,
\[
\varphi(z)=4(z-e_1)(z-e_2)(z-e_3)=4z^3-g_2z-g_3
\]
sein soll,
\begin{equation}
g_2=12\mu^2A,\qquad g_3=-4\mu^3B.
\tag{6}
\end{equation}
Hierin sind
\begin{equation}
\begin{aligned}
A&=a_2^2-3a_1a_3+12a_0a_4,\\
B&=27a_1^2a_4+27a_0a_3^2+2a_2^3-72a_0a_2a_4-9a_1a_2a_3,
\end{aligned}
\tag{7}
\end{equation}
die erste und zweite Invariante der biquadratischen Form $f(x)$.

Die logarithmische Differentiation der beiden ersten Gleichungen (4) ergibt
\[
\frac{dz^2}{(z-e_1)(z-e_2)}=\frac{dx^2(x_1x_2)(x_1x_3)}{(xx_1)^2(xx_2)(xx_3)},
\]
und mit Hilfe der letzten Gleichung (4) nach (5):
\begin{equation}
\frac{dz}{\sqrt{4z^3-g_2z-g_3}}=\frac{dx}{\sqrt{12\mu f(x)}}.
\tag{8}
\end{equation}

Wollen wir diese Resultate auf den Fall anwenden, wo $f(x)$ die Normalform
\[
f(x)=x(1-x)(1-\varkappa^2x)=x-x^2(1+\varkappa^2)+\varkappa^2x^3
\]
hat, und der Punkt $z=\infty$ dem Punkte $x=0$ entspricht, so lassen wir $a_0$ in Null, $x_4$ in Unendlich übergehen, aber das Produkt $a_0x_4$ in einen endlichen Wert, den wir $=1$ annehmen können.

Es wird dann
\[
a_0x_4=1,\qquad a_0=0,\qquad a_1=\varkappa^2,
\]
\[
a_2=-(1+\varkappa^2),\qquad a_3=1,\qquad a_4=0,
\]
und folglich
\[
A=(1+\varkappa^2)^2-3\varkappa^2=1-\varkappa^2-\varkappa^4=1-\varkappa^2\varkappa'^2,
\]
\[
\begin{aligned}
B&=-(1+\varkappa^2)\{2(1+\varkappa^2)^2-9\varkappa^2\}\\
&=-(1+\varkappa^2)(2-\varkappa^2)(1-2\varkappa^2)\\
&=(1+\varkappa^2)(1+\varkappa'^2)(\varkappa^2-\varkappa'^2)\\
&=(2+\varkappa^2\varkappa'^2)(\varkappa^2-\varkappa'^2),
\end{aligned}
\]
wenn $\varkappa'^2=1-\varkappa^2$ gesetzt ist.

Es wird also, wenn wir noch die Diskriminante
\[
\Delta=16\cdot 27\,\mu^6(4A^3-B^2)
\]
beifügen,
\begin{equation}
\begin{aligned}
g_2&=12\mu^2(1-\varkappa^2\varkappa'^2),\\
g_3&=-4\mu^3(2+\varkappa^2\varkappa'^2)(\varkappa^2-\varkappa'^2),\\
\Delta&=g_2^3-27g_3^2=27^2\cdot16\,\mu^6\varkappa^4\varkappa'^4,
\end{aligned}
\tag{9}
\end{equation}
und man erhält aus (6) die Transformation
\begin{equation}
\frac{dz}{\sqrt{4z^3-g_2z-g_3}}=
\frac{dx}{\sqrt{12\mu x(1-x)(1-\varkappa^2x)}};
\tag{10}
\end{equation}
um die Substitution zu finden, setzen wir $x_1=0$, $x_2=1$, $x_3=1:\varkappa^2$, $x_4=\infty$, $a_0x_4=1$ und erhalten
\[
a_0U=1,\qquad a_0V=-\frac1{\varkappa^2},\qquad a_0W=\frac{\varkappa'^2}{\varkappa^2},
\]
folglich
\[
e_1=\mu\frac{\varkappa^2-2}{\varkappa^2},\qquad
 e_2=\mu\frac{1-2\varkappa^2}{\varkappa^2},\qquad
 e_3=\mu\frac{1+\varkappa^2}{\varkappa^2},
\]
\begin{equation}
z=\frac{3\mu}{\varkappa^2}\left(\frac{\varkappa^2+1}{3}-\frac1x\right).
\tag{11}
\end{equation}

Bei dieser Transformation des elliptischen Differentials in der Weierstrassschen Normalform wird die Zerlegung der Funktion $f(x)$ in ihre linearen Faktoren, also die Auflösung der biquadratischen Gleichung $f(x)=0$, vorausgesetzt. Eine andere, freilich nicht lineare Transformation, die ohne diese Voraussetzung den gleichen Zweck erreicht, hat Hermite gegeben (Crelles Journal, Bd.~52). Um sie darzustellen, benutzen wir die homogene Form des elliptischen Differentials
\begin{equation}
\frac{y\,dx-x\,dy}{\sqrt{f(x,y)}},
\tag{12}
\end{equation}
worin
\begin{equation}
f(x,y)=a_0x^4+a_1x^3y+a_2x^2y^2+a_3xy^3+a_4y^4.
\tag{13}
\end{equation}
Diese biquadratische Form hat außer den beiden Invarianten $A,B$ noch zwei Kovarianten:
\begin{equation}
H=\frac13\left[\frac{\partial^2 f}{\partial x^2}\frac{\partial^2 f}{\partial y^2}-\left(\frac{\partial^2 f}{\partial x\partial y}\right)^2\right],\qquad
T=\frac1{12}\left(\frac{\partial f}{\partial x}\frac{\partial H}{\partial y}-\frac{\partial f}{\partial y}\frac{\partial H}{\partial x}\right),
\tag{14}
\end{equation}
wo $H$ und $T$ Formen vierten und sechsten Grades sind, deren Koeffizienten sich rational und mit ganzzahligen Zahlenkoeffizienten aus den Koeffizienten von $f$ zusammensetzen. Zwischen diesen Formen besteht dann noch die Relation
\begin{equation}
H^3-48AHf^2-64Bf^3=-27T^2
\tag{15}
\end{equation}
(Bd.~I, \S\S~70, 72).

Es ist aber nach dem Eulerschen Satz über homogene Funktionen
\[
4f=\frac{\partial f}{\partial x}x+\frac{\partial f}{\partial y}y,\qquad
df=\frac{\partial f}{\partial x}dx+\frac{\partial f}{\partial y}dy,
\]
\[
4H=\frac{\partial H}{\partial x}x+\frac{\partial H}{\partial y}y,\qquad
dH=\frac{\partial H}{\partial x}dx+\frac{\partial H}{\partial y}dy,
\]
woraus
\[
f\,dH-H\,df=-3T(y\,dx-x\,dy),
\]
und wenn also nach (15)
\begin{equation}
3\sqrt3\,T=\sqrt{-H^3+48Af^2H+64Bf^3}
\tag{16}
\end{equation}
gesetzt wird:
\[
\frac{y\,dx-x\,dy}{\sqrt f}=\frac{-\sqrt3(f\,dH-H\,df)}{\sqrt{-(H^3-48Af^2H-64Bf^3)f}}.
\]
Macht man nun die Substitution
\begin{equation}
z=-\frac{\mu H}{4f}
\tag{17}
\end{equation}
mit einem unbestimmten Faktor $\mu$, so ergibt sich
\begin{equation}
\frac{y\,dx-x\,dy}{\sqrt{3\mu f}}=\frac{dz}{\sqrt{4z^3-g_2z-g_3}},
\tag{18}
\end{equation}
wenn wie früher
\begin{equation}
g_2=12\mu^2A,\qquad g_3=-4\mu^3B.
\tag{19}
\end{equation}

Man kann diese Transformation auf ein Differential anwenden, das schon die Normalform hat. Setzt man
\[
f(x,y)=4x^3y-g_2xy^3-g_3y^4,
\]
so hat man
\[
\begin{array}{ccccc}
a_0&a_1&a_2&a_3&a_4\\
0&4&0&-g_2&-g_3
\end{array}
\]
zu ersetzen, und es ergibt sich aus (7)
\[
A=12g_2,
\qquad B=-27\cdot16\,g_3.
\]
Wenn man also $\mu=\frac1{12}$ setzt, so gehen die Gleichungen (19) in die Identitäten $g_2=g_2$, $g_3=g_3$ über.

Wenn man dann $y=1$ setzt, so ergibt die Transformation (18)
\begin{equation}
\frac{2dx}{\sqrt{4x^3-g_2x-g_3}}=\frac{dz}{\sqrt{4z^3-g_2z-g_3}}.
\tag{20}
\end{equation}
Diese Transformation wird nach (16) durch
\[
z=-\frac{H}{48f}
\]
vermittelt. Es ergibt sich aber aus (14):
\[
H=-48\left[\left(x^2+\frac14g_2\right)^2+2g_3x\right],
\]
und folglich erhalten wir
\begin{equation}
z=\frac{\left(x^2+\frac14g_2\right)^2+2g_3x}{4x^3-g_2x-g_3}.
\tag{21}
\end{equation}
Weiter ergibt sich noch aus (14)
\[
T=64x^6-80g_2x^4-320g_3x^3-20g_2^2x^2-16g_2g_3x+g_2^3-32g_3^2
\]
und dann nach (16)
\begin{equation}
4\left(\sqrt{4x^3-g_2x-g_3}\right)^3\sqrt{4z^3-g_2z-g_3}=T.
\tag{22}
\end{equation}
Durch (20), (21), (22) ist die Multiplikation des elliptischen Differentials mit $2$ geleistet. Es ist dadurch nicht nur $z$ als eindeutige (rationale) Funktion von $x$ dargestellt, sondern auch die eine Quadratwurzel eindeutig durch die andere, d.~h. es ist einem Punkte $x$ eindeutig ein Punkt $z$ zugeordnet.

\section*{\S\,6. Elliptische Kurven.}

Jede algebraische Abhängigkeit zwischen zwei Veränderlichen $x,y$ wird ausgedrückt durch eine Gleichung der Form
\begin{equation}
F(x,y)=0,
\tag{1}
\end{equation}
worin $F(x,y)$ eine ganze Funktion der beiden Veränderlichen $x,y$ bezeichnet. Betrachtet man $x$ und $y$ als Cartesische Koordinaten eines Punktes in der Ebene, so ist (1) die Gleichung einer Kurve $n$ten Grades, wenn $F$ in bezug auf $x$ und $y$ zusammengenommen von der $n$ten Dimension ist. Diese Kurve ist dann das geometrische Bild der algebraischen Abhängigkeit, wobei indessen auch imaginäre Punkte mit berücksichtigt werden müssen.

Ist dann $\Phi(x,y)$ eine ganze oder gebrochene rationale Funktion von $x$ und $y$, und $y$ von $x$ durch die Gleichung (1) abhängig, so sind
\begin{equation}
\Phi(x,y)\,dx,\qquad \int\Phi(x,y)\,dx
\tag{2}
\end{equation}
die zu der Kurve $F$ gehörigen algebraischen Differentiale und Integrale.

Beispielsweise ist, wenn $f(x)$ eine Funktion dritten oder vierten Grades von $x$ ist,
\[
y^2-f(x)=0
\]
die Gleichung einer Kurve dritten oder vierten Grades, zu der die mit der Irrationalität $\sqrt{f(x)}$ behafteten elliptischen Differentiale und Integrale gehören.

Wir wollen hier alle Kurven, deren Differentiale und Integrale auf elliptische reduzierbar sind, elliptische Kurven nennen. Wie das Beispiel zeigt, kommen darunter Kurven dritten und vierten Grades vor.

Kegelschnitte gehören nicht zu den elliptischen Kurven, weil, wenn zwischen $x$ und $y$ eine Gleichung zweiten Grades besteht, $x$ und $y$ als rationale Funktionen eines Parameters $t$ dargestellt werden können, wodurch das algebraische Differential auf ein rationales nach $t$ zurückgeführt werden kann. Solche Kurven heißen rationale Kurven.

Wir werden sehen, daß alle Kurven dritten Grades ohne Doppel- oder Rückkehrpunkt zu den elliptischen gehören. Kurven höheren Grades können nur dann dazu gehören, wenn sie eine gewisse Anzahl singulärer Punkte haben. Dies ist ein fundamentales Kapitel in der allgemeinen Theorie der algebraischen Funktionen, das nicht in den Plan dieses Werkes gehört.\footnote{Die wichtigsten diesen Gegenstand betreffenden Arbeiten sind: Riemann, Theorie der Abelschen Funktionen [Crelles Journal, Bd.~54 (1857)]. Mathematische Werke, 2.~Aufl., S.~88, 467. Nachträge, herausgegeben von Noether und Wirtinger (1902). Aronhold, Monatsberichte der Berliner Akademie vom 25.~April 1861. Clebsch, Über die Anwendung der Abelschen Funktionen in der Geometrie (Crelles Journal, Bd.~63). Clebsch, Über diejenigen ebenen Kurven, deren Koordinaten rationale Funktionen eines Parameters sind (ibid., Bd.~64). Clebsch, Über diejenigen Kurven, deren Koordinaten sich als elliptische Funktionen eines Parameters darstellen lassen (ibid., Bd.~64). Brill, Über diejenigen Kurven, deren Koordinaten sich als hyperelliptische Funktionen eines Parameters darstellen lassen (ibid., Bd.~65). Clebsch und Gordan, Theorie der Abelschen Funktionen (Teubner, Leipzig 1866).}

Für die Untersuchung algebraischer Differentiale von diesem Gesichtspunkte ist die Einführung homogener Variablen zweckmäßig. Wir setzen $x=x_1:x_3$, $y=x_2:x_3$, $x_3^nF(x,y)=f(x_1,x_2,x_3)$; dann ist $f(x_1,x_2,x_3)$ eine homogene Funktion $n$ter Ordnung der drei Veränderlichen $x_1,x_2,x_3$ und
\begin{equation}
f(x_1,x_2,x_3)=0
\tag{3}
\end{equation}
die Gleichung einer Kurve $n$ter Ordnung in homogenen Koordinaten. Es wird
\[
dx=\frac{x_3\,dx_1-x_1\,dx_3}{x_3^2}
\]
und
\begin{equation}
d\Omega=\Phi(x,y)dx=\frac1{x_3^2}\Phi\left(\frac{x_1}{x_3},\frac{x_2}{x_3}\right)(x_3\,dx_1-x_1\,dx_3).
\tag{4}
\end{equation}
Nun ist aber nach dem Eulerschen Satz über homogene Funktionen, wenn wir mit $f_1,f_2,f_3$ die partiellen Ableitungen von $f$ nach $x_1,x_2,x_3$ bezeichnen,
\[
f_1x_1+f_2x_2+f_3x_3=0,
\qquad
f_1dx_1+f_2dx_2+f_3dx_3=0,
\]
daher, wenn $\rho$ einen Proportionalitätsfaktor bedeutet,
\[
f_1=\rho(x_2\,dx_3-x_3\,dx_2),\qquad
f_2=\rho(x_3\,dx_1-x_1\,dx_3),\qquad
f_3=\rho(x_1\,dx_2-x_2\,dx_1),
\]
und wenn man mit $c_1,c_2,c_3$ ganz willkürliche Größen, z.~B. Konstanten, bezeichnet:
\[
c_1f_1+c_2f_2+c_3f_3=\rho\sum\pm c_1x_2\,dx_3,
\]
wenn $\sum\pm c_1x_2\,dx_3$ in üblicher Weise die Determinante
\[
\begin{vmatrix}
c_1&c_2&c_3\\
x_1&x_2&x_3\\
dx_1&dx_2&dx_3
\end{vmatrix}
\]
bedeutet. Es folgt also:
\[
x_3\,dx_1-x_1\,dx_3=\frac{f_2\sum\pm c_1x_2\,dx_3}{c_1f_1+c_2f_2+c_3f_3},
\]
und wenn man
\[
\frac{f_2}{x_3^2}\Phi\left(\frac{x_1}{x_3},\frac{x_2}{x_3}\right)=\Psi
\]
setzt, so ergibt sich aus (4) der Ausdruck für das allgemeinste zu der Kurve $f$ gehörige algebraische Differential:
\begin{equation}
d\Omega=\frac{\Psi\sum\pm c_1x_2\,dx_3}{c_1f_1+c_2f_2+c_3f_3},
\tag{5}
\end{equation}
worin $\Psi$ eine ganze oder gebrochene homogene Funktion der $(n-3)$ten Ordnung ist. Die willkürlichen Größen $c_1,c_2,c_3$ kommen nur scheinbar in diesem Ausdruck vor. In Wirklichkeit ist er davon ganz unabhängig.

Wir betrachten wieder den Fall $n=3$. Dann ist der einfachste Fall der, daß $\Psi$ eine Konstante ist, und (5) geht dann in das elliptische Differential erster Gattung über:
\begin{equation}
du=\frac{\sum\pm c_1x_2\,dx_3}{c_1f_1+c_2f_2+c_3f_3}.
\tag{6}
\end{equation}
Dieses wollen wir nun mit Hilfe der Kovarianten der ternären Form dritten Grades auf die Weierstrasssche Form transformieren.

Wir haben Bd.~II, \S~107 die fundamentalen Kovarianten der kubischen Form kennen gelernt. Es waren dies außer $f$ selbst:
\[
\Delta=\frac1{6^3}\begin{vmatrix}
f_{11}&f_{12}&f_{13}\\
f_{21}&f_{22}&f_{23}\\
f_{31}&f_{32}&f_{33}
\end{vmatrix},\qquad
J=\frac1{36}\begin{vmatrix}
f_{11}&f_{12}&f_{13}&\Delta_1\\
f_{21}&f_{22}&f_{23}&\Delta_2\\
f_{31}&f_{32}&f_{33}&\Delta_3\\
\Delta_1&\Delta_2&\Delta_3&0
\end{vmatrix},
\]
\[
K=\frac19\begin{vmatrix}
f_1&\Delta_1&J_1\\
f_2&\Delta_2&J_2\\
f_3&\Delta_3&J_3
\end{vmatrix},
\]
wobei die Indizes $1,2,3$ die Differentiation nach den Variablen $x_1,x_2,x_3$ bedeuten.

Es folgt aus dem Multiplikationssatz der Determinanten, da $f$ und $df=0$ sind:
\[
\begin{vmatrix}
c_1&c_2&c_3\\
x_1&x_2&x_3\\
dx_1&dx_2&dx_3
\end{vmatrix}
\begin{vmatrix}
f_1&f_2&f_3\\
\Delta_1&\Delta_2&\Delta_3\\
J_1&J_2&J_3
\end{vmatrix}
=3\sum c_if_i(\Delta\,dJ-2J\,d\Delta),
\]
und daraus
\begin{equation}
du=3\frac{\Delta\,dJ-2J\,d\Delta}{K}.
\tag{7}
\end{equation}

Mit Hilfe der Kovarianten haben wir die Form $f(x)$ auf die kanonische Form
\[
\varphi(y)=y_1^3+y_2^3+y_3^3+6my_1y_2y_3
\]
reduziert. Wir haben, wenn $r$ die Substitutionsdeterminante bedeutet:
\[
P=\frac{(2+m^3)m^3f^2+(2-5m^3)mr^2f\Delta+3m^2r^4\Delta^2-r^6J}{(1+8m^3)^2},
\]
\[
Q=\frac{(1+2m^3)f-6mr^2\Delta}{1+8m^3},\qquad
R=\frac{m^2f+r^2\Delta}{1+8m^3}
\]
gesetzt und erhielten $y_1^3,y_2^3,y_3^3$ als Wurzeln der kubischen Gleichung
\[
u^3-Qu^2+Pu-R^3=0,
\]
und die Diskriminante dieser kubischen Gleichung ist
\[
D_1=\frac{r^{18}K^2}{(1+8m^3)^6}.
\]
Drückt man diese Diskriminante durch $P,Q,R$ aus, so erhält man $K^2$ rational durch $f,\Delta,J$ dargestellt. Wir haben an der erwähnten Stelle der Algebra diesen Ausdruck nicht explizite angegeben. Jetzt müssen wir ihn aber bilden, wenn auch nur unter der Voraussetzung $f=0$, d.~h. nur für die Punkte der Kurve. Es ist aber [nach Bd.~I, \S~50, (10)]
\[
D_1=P^2Q^2+18PQR^3-4P^3-4Q^3R^3-27R^6.
\]
Darin hat man zu setzen:
\[
P=\frac{3m^2r^4\Delta^2-r^6J}{(1+8m^3)^2},\qquad
Q=\frac{-6mr^2\Delta}{1+8m^3},\qquad
R=\frac{r^2\Delta}{1+8m^3},
\]
und man erhält durch eine nicht schwierige Rechnung, wenn man mit $S$ und $T$ die beiden Invarianten der Kurve dritter Ordnung bezeichnet (Bd.~II, \S~108),
\[
(1+8m^3)^6D_1=r^{18}(4J^3+108SJ\Delta^4-27T\Delta^6),
\]
und folglich
\begin{equation}
K^2=4J^3+108SJ\Delta^4-27T\Delta^6.
\tag{8}
\end{equation}
Diese Gleichung ist aber nicht identisch, sondern nur unter der Voraussetzung $f=0$ befriedigt. Der vollständige Ausdruck von $K^2$ durch $J,\Delta,f$ wird sehr viel komplizierter.

Setzt man
\begin{equation}
\frac{J}{\Delta^2}=3z,\qquad dz=\frac13\frac{\Delta\,dJ-2J\,d\Delta}{\Delta^3},
\tag{9}
\end{equation}
so ergibt sich aus (8)
\[
K^2=27\Delta^6(4z^3+12Sz-T)
\]
und aus (7)
\begin{equation}
du=\frac1{\sqrt3}\frac{dz}{\sqrt{4z^3+12Sz-T}},
\tag{10}
\end{equation}
und dies ist die Weierstrasssche Normalform für $g_2=-12S$, $g_3=T$.

\end{document}
